\section{Background and Related Works} \label{sec:back}
This section provides a brief literature review, extracted mainly from Scopus, on the current state of the art of the main topics discussed by this study.

\subsection{Network Anomaly Detection (NAD) in IoT}
% TITLE-ABS-KEY ( ( "intrusion detection systems" OR "anomaly detection" ) AND ( "network security" OR "cybersecurity" ) AND ( "network traffic" OR "flow" ) ) AND PUBYEAR > 2023 AND PUBYEAR < 2027 AND ( LIMIT-TO ( DOCTYPE , "re" ) ) AND ( LIMIT-TO ( LANGUAGE , "English" ) )

NAD refers to the analysis of the network traffic to detect abnormal behaviors, which can be either legitimate changes or attacks. In the IoT context, detecting these events is crucial to ensure the integrity of decision systems, where failure or downtime can result in risks and costs \cite{liso:2024}. 

Traditionally, as in \cite{ge:2019}, the subject of this analysis was network packets, inspected through a deep search of their content or headers. With the increase in network traffic, this approach proved to be neither scalable nor computationally effective \cite{wang:2025,miguel:2026}. A well-established approach to reducing data volume while preserving relevant behavioral information is the use of network flows. Network flows summarize communication sessions between hosts by aggregating packets that share common header characteristics, such as source and destination addresses, ports, and protocols \cite{miguel:2026}. This makes flow-based analysis suitable for scalable monitoring and threat detection in high-speed networks.

The network flow analysis can be conducted using signature-based or anomaly-based methods \cite{elnakib:2023}. In signature-based analysis, a knowledge database stores all the already mapped attack signatures. The flow patterns are compared with the stored patterns, and if a match occurs, it is marked as anomalous or malicious. Usually, rule-based models have better accuracy scores. The disadvantages of this approach are the rapid obsolescence of the database as new types of attacks emerge, failing to detect zero-day attacks, where the attack patterns are unknown or significantly different from previously observed threats. To address this, anomaly-based methods are better suited, since they focus on learning the traffic normal behavior \cite{anis:2025, elnakib:2023}. All observations that differ from the learned behavior are classified as an anomaly. Although they can deal better with unknown patterns, they lack accuracy, raising more false positives when compared with signature-based methods. In the literature, implementations usually adopt the anomaly-based solution, although, as implemented by \cite{hussain:2024}, a hybrid approach proved to be the best way to cover both known and unknown anomalies.

% TODO: pegar as referencias dos datasets
The uprising of data traffic in the IoT network brings challenges regarding the best method to learn the normal traffic behavior. That is why deep learning (DL) methods are in the spotlight of NAD literature \cite{anis:2025, tsimenidis:2022}. They can learn patterns from a great amount of historical data and classify traffic with a high accuracy score \cite{barnawi:2023}. These techniques can be broadly categorized into supervised and unsupervised approaches. Supervised methods rely on labeled data and often require detailed contextual information, including prior knowledge of attack characteristics \cite{sharma:2025, miguel:2026}. Most supervised methods implementations use public datasets to acquire labeled data, such as CICIDS2017, NSL–KDD and KDDCup99. In \cite{elnakib:2023, toupas:2019, ge:2019}, although a high accuracy was obtained by these models, this strategy resulted in a less suitable approach for detecting zero-day attacks \cite{miguel:2026}. In unsupervised techniques, data is unlabeled, and no ground-truth is used. Learning comes from hidden data patterns, identifying structures and relationships to group common behaviors. Outliers or deviations are detected when there are no significant shared similarities with other groups. All the knowledge required for pattern detection is extracted from data, with no prior information needed. As demonstrated by \cite{hussain:2024} approach, this makes unsupervised methods more suitable for detecting new anomaly signatures \cite{sharma:2025, miguel:2026}. Besides that, few studies chose unsupervised approaches \cite{hussain:2024}.

% TITLE-ABS-KEY ( ( "intrusion detection systems" OR "anomaly detection" ) AND ( "network security" OR "cybersecurity" ) AND ( "internet of things" OR "IoT" ) ) AND PUBYEAR > 2023 AND PUBYEAR < 2027 AND ( LIMIT-TO ( DOCTYPE , "re" ) ) AND ( LIMIT-TO ( LANGUAGE , "English" ) )

Regarding the DL learning models for NAD, the predominant ones in the IoT literature are AutoEncoder, Multi-Layer Perceptron, Long Short-Term Memory, Deep Neural Network, and Convolutional Neural Network \cite{elnakib:2023,kalpani:2026}. In \cite{kalpani:2026}, the trade-offs of AI techniques were discussed; concerns about limited computational capabilities must be addressed in IoT networks, specially in edge computing. Despite many implementations of DL techniques being present in the literature \cite{elnakib:2023, toupas:2019, ge:2019, hussain:2020}, few of them assess the computational efficiency and scalability of their proposal. The majority of the works assessed detection metrics, such as accuracy, true- and false-positives rates, and ROC curves. There is a lack of energy, memory, and time complexity assessment \cite{di:2021}.

Also, the use of public datasets limits the approaches to specific mapped attacks, specific context, and outdated traffic volume and features \cite{sharafaldin:2018, di:2021}. True deployments of the proposals are not tested to evaluate real load compliance. 

\subsection{Weightless Neural Network in IoT}
Inspired by the decode processing in the dendritic trees of biological neurons, WNN presents RAM nodes as an alternative to weighted-sum-and-threshold neurons in DNN \cite{aleksander:2009}. RAM nodes use lookup tables (LUTs) to convert entries in Boolean functions. This approach reduces the number of multiply-accumulate operations, making it a cost-efficient solution.

Its first implementation was Wilkie, Stonham, and Aleksander’s Recognition Device (WiSARD) in 1984. WiSARD uses a set of discriminators composed of RAM neurons to learn binary patterns associated with each class \cite{di:2020}. Although WiSARD has low computational load and low inference latency requirements, it requires a considerable amount of memory to store RAM nodes, with direct implications for the size of the input \cite{lello:2020}. These drawbacks were addressed in \cite{ferreira:2019}, using hash tables to decrease the architecture's memory consumption, and in \cite{miranda:2022}, using only logic functions to avoid memories and arithmetic circuits. Among the most significant results in memory consumption reduction is the Bloom filter-based strategy, presented by \cite{santiago:2020, susskind:2022}, where probabilistic membership queries are used instead of deterministic RAM lookup.

Though memory limitation was overcome by new studies, the earlier implementations of WNN still fail to surpass DNNs in accuracy. However, demonstrations of WNN superiority over quantized DNNs and the possibility to implement it in field-programmable gate array (FPGA) and ASICs make it a promising alternative for IoT edge environments \cite{susskind:2023}. In \cite{ferreira:2019, susskind:2022}, the implementations of WiSARD in FPGA demonstrated its adequacy for resource-restricted devices. In \cite{nag:2024}, the study compares two LUT-based neural networks performance in an FPGA, proving their cost-efficiency in IoT. That supports this study proposal in applying WNN in IoT networks.

%% TODO: falar de encoding e bleach
Another challenge in the WiSARD model is the encoding of the inputs; the wrong strategy can cause loss of information, which directly impact in the classification results \cite{kappaun:2016}. That is why special attention was given to this subject. Following the \cite{kappaun:2016} study, the encoding approach must consider Hamming distance \cite{hamming:1950} to preserve data similarity and limit the impact of small numerical variations. This study's experiments showed that, in general, thermometer encoding achieved the best accuracy score against other techniques, given its characteristic of being robust to noise. Also, bleaching technique and random RAM mapping presents better representation for anomaly detection in \cite{susskind:2022, aleksander:2009}.

Few studies have addressed WiSARD as a solution for IDS, especially considering the IoT context. In both works of \cite{di:2020, di:2021}, an experimental evaluation of WiSARD using public IDS datasets is conducted against baseline neural networks. WiSARD presented high performance for most of the tested datasets. When considering average accuracy and time complexity, WiSARD delivered the best cost-benefit equilibrium, as stressed by the authors.

%% Final da seção
Unlike previous studies, this work proposes the use of WNN for NAD in the IoT context as a security solution for resource-restricted networks. To the best of our knowledge, this work is the first to use WNN in IoT network security. The core contributions are the proposal of an encoding scheme for network flow entries in WNN, calibration strategies for hyperparameters, and experimentation design to assess computational and performance requirements using public datasets and real traffic environments.